\begin{figure}[t]
  \noindent\textbf{\footnotesize (a) Real source bug (bcc biosnoop.bpf.c)}
  \begin{lstlisting}[style=bpfix,language=C,
      linebackgroundcolor={\ifnum\value{lstnumber}=2 \color{lstbad!40}\fi}]
// rq is a raw kernel pointer from a kprobe
event.sector = rq->__sector;     // direct deref: rq is a scalar
event.len    = rq->__data_len;
  \end{lstlisting}
  \begin{lstlisting}[style=bpflog]
verifier: R1 invalid mem access 'scalar'
  \end{lstlisting}

  \vspace{3pt}
  \noindent\textbf{\footnotesize (b) Compiler artifact (cilium nat.h)}
  \begin{lstlisting}[style=bpfix,language=C,
      linebackgroundcolor={\ifnum\value{lstnumber}=3 \color{lstgood!40}\fi}]
// correct source; lowering drops the packet-pointer provenance
tuple.nexthdr = iphdr.nexthdr;
hdrlen = ipv6_hdrlen_offset(ctx, inner_l3_off, &tuple.nexthdr, &fraginfo);
  \end{lstlisting}
  \begin{lstlisting}[style=bpflog]
verifier: R7 invalid mem access 'scalar'
  \end{lstlisting}

  \caption{Two programs that reject with the same message, each shown with the verifier's output. (a) dereferences a kernel pointer directly, a real source bug fixed with a CO-RE read. (b) is correct source that the compiler lowered into an untracked packet-pointer path, fixed by forcing the value onto the stack. The two fixes are opposite. Of the 28 cases with this message, 23 are real source bugs, four are compiler artifacts, and one is an environment issue.}
  \label{fig:packet}
\end{figure}
